\begin{table}[h]
\centering\small
\begin{threeparttable}
\caption{The five declared matching criteria.}
\label{tab:si-criteria}
\begin{tabular}{lp{7.2cm}l}
\toprule
criterion & what counts as the same structure & strictness \\
\midrule
\texttt{canonical} & RDKit canonical SMILES, equality of strings & strictest \\
\texttt{inchikey} & the full InChIKey, including the stereochemistry block & strict \\
\texttt{inchi\_no\_stereo} & the first block of the InChIKey. This is the skeleton hash, so it drops the protonation and isotope layers along with stereochemistry & medium \\
\texttt{tanimoto1} & Tanimoto similarity of one on 1024-bit Morgan fingerprints & loosest \\
\texttt{inchikey\_tautomer} & canonical tautomer on both sides, then a canonical SMILES without stereochemistry & the default \\
\bottomrule
\end{tabular}
\begin{tablenotes}[flushleft]\footnotesize
\item The references carry no stereochemistry and do not distinguish tautomers, and standard InChI normalises only a subset of tautomers, which is why the tautomer-aware key is the default rather than the full InChIKey.
\end{tablenotes}
\end{threeparttable}
\end{table}
